%% toc: P15 158 158.3,158.4
% Proofs for the integration supplements: monotone limits, signed series, convergence in measure, Jensen, moment bounds, inclusion-exclusion and Bonferroni (Durrett, secs. 1.5-1.6).
\section{Integration supplements: proofs}
\label{sec:annex-p15-w02c-p15-005-integration-supplements}

Use the definitions of
Chapter~\ref{sec:p15-w02c-p15-005-integration-supplements} and
Chapter~\ref{sec:p15-integration-expectation}. Two facts are used
repeatedly. First, if two extended real measurable functions agree a.e.,
then so do their positive parts and their negative parts, so by
\cref{thm:p15-nonnegative-integral-laws}~(2) the two part integrals
agree and the signed integrals exist together and agree. Second, a
nonnegative measurable function with finite integral is finite a.e., by
\cref{thm:p15-nonnegative-integral-laws}~(1).

\begin{annexitem}{annex:p15-w02c-p15-005-monotone-lower}{thm:p15-w02c-p15-005-monotone-lower}{Shifting by the negative part of the first term}
\begin{proof}
Let \(N\) be a measurable null set off which \(g_n\leq g_{n+1}\) for
every \(n\) and \(g_n\to g\). Since \(\int g_1^-\,d\mu<\infty\), the
set \(N'=\{g_1^-=\infty\}\) is null. Put \(G=(N\cup N')^c\). On \(G\),
\(g_1\leq g_n\leq g\), so
\[
 g_n^-\leq g_1^-<\infty,\qquad g^-\leq g_1^-<\infty .
\]
Hence \(\int g_n^-\,d\mu\) and \(\int g^-\,d\mu\) are at most
\(\int g_1^-\,d\mu<\infty\) (the inequalities hold a.e.), so every
signed integral in the statement exists and lies in
\((-\infty,\infty]\). By the first preliminary fact, replacing each
\(g_n\) and \(g\) by its product with \(\mathbf1_G\) changes no
integral; so assume the monotonicity, the convergence and the
finiteness of \(g_1^-\) hold at every point, with all functions zero on
\(G^c\).

Define
\[
 h_n\defeq g_n^++(g_1^--g_n^-),\qquad h\defeq g^++(g_1^--g^-).
\]
Both summands are nonnegative and measurable, and the differences are
finite real numbers. At each point, \(h_n=g_n+g_1^-\) and
\(h=g+g_1^-\), where the right sides are formed with the finite number
\(g_1^-\); this holds also where \(g_n=\infty\), since then
\(g_n^-=0\). As \(g_n\) increases to \(g\) and \(g_1^-\) is a fixed
finite number at each point, \(h_n\) increases to \(h\). Monotone
convergence (\cref{thm:p15-monotone-integral}) gives
\(\int h_n\,d\mu\uparrow\int h\,d\mu\).

Additivity of the nonnegative integral
(\cref{thm:p15-nonnegative-integral-laws}) gives
\(\int(g_1^--g_n^-)\,d\mu+\int g_n^-\,d\mu=\int g_1^-\,d\mu\). All three
terms are finite, so
\(\int(g_1^--g_n^-)\,d\mu=\int g_1^-\,d\mu-\int g_n^-\,d\mu\). Again by
additivity,
\[
 \int h_n\,d\mu=\int g_n^+\,d\mu-\int g_n^-\,d\mu+\int g_1^-\,d\mu
 =\int g_n\,d\mu+\int g_1^-\,d\mu,
\]
and the same computation gives
\(\int h\,d\mu=\int g\,d\mu+\int g_1^-\,d\mu\). Subtracting the finite
number \(\int g_1^-\,d\mu\) from the increasing convergence
\(\int h_n\,d\mu\uparrow\int h\,d\mu\) gives
\(\int g_n\,d\mu\uparrow\int g\,d\mu\), including the case where the
limit is \(\infty\).

On a probability space take \(\mu=\mathbb P\), \(g_n=X_n\) and
\(g=X\); the hypothesis \(\mathbb EX_1^-<\infty\) is
\(\int g_1^-\,d\mathbb P<\infty\), and the conclusion is
\(\mathbb EX_n\uparrow\mathbb EX\).
\end{proof}
\end{annexitem}

\begin{annexitem}{annex:p15-w02c-p15-005-signed-series}{thm:p15-w02c-p15-005-signed-series}{Domination by the series of absolute values}
\begin{proof}
The function \(G=\sum_n|f_n|\) is a nonnegative measurable function, a
countable sum of nonnegative measurable functions. By the series
identity of \cref{thm:p15-nonnegative-integral-laws},
\(\int G\,d\mu=\sum_n\int|f_n|\,d\mu<\infty\), so \(G\) is integrable
and the set \(E=\{G<\infty\}\) has null complement. At each
\(\omega\in E\) the real series \(\sum_nf_n(\omega)\) converges
absolutely, so its partial sums \(S_N=\sum_{n\leq N}f_n\) converge
there.

The functions \(S_N\mathbf1_E\) are real measurable and converge at
every point to \(F\), which equals \(\sum_nf_n\) on \(E\) and zero on
\(E^c\). So \(F\) is measurable, as a pointwise limit of real measurable
functions. Moreover
\(|S_N\mathbf1_E|\leq\sum_{n\leq N}|f_n|\leq G\) everywhere, and \(G\)
is integrable. Dominated convergence
(\cref{thm:p15-dominated-bounded-convergence}~(1)) shows that \(F\) is
integrable and \(\int S_N\mathbf1_E\,d\mu\to\int F\,d\mu\). Each
\(f_n\) is integrable, since \(|f_n|\leq G\), and \(f_n\mathbf1_E=f_n\)
a.e.; by linearity (\cref{thm:p15-signed-integral-laws}),
\[
 \int S_N\mathbf1_E\,d\mu=\sum_{n\leq N}\int f_n\mathbf1_E\,d\mu
 =\sum_{n\leq N}\int f_n\,d\mu .
\]
Hence \(\sum_n\int f_n\,d\mu\) converges and equals \(\int F\,d\mu\).
It converges absolutely because
\(\left|\int f_n\,d\mu\right|\leq\int|f_n|\,d\mu\) and the latter terms
have a finite sum.

For extended real integrable \(f_n\), each \(f_n\) is finite a.e., and
its real-valued representative \(\tilde f_n\) of
\cref{def:p15-integrable-functions} equals \(f_n\) a.e. Then
\(\int\tilde f_n\,d\mu=\int f_n\,d\mu\) and
\(\int|\tilde f_n|\,d\mu=\int|f_n|\,d\mu\), so the first part applies to
\((\tilde f_n)\). The countably many exceptional sets have null union,
so the resulting \(F\) agrees a.e.\ with \(\sum_nf_n\) wherever the
latter converges absolutely.
\end{proof}
\end{annexitem}

\begin{annexitem}{annex:p15-w02c-p15-005-indefinite-continuity}{thm:p15-w02c-p15-005-indefinite-continuity}{Sequential continuity by dominated convergence}
\begin{proof}
Let \(x\in[a,b]\) and let \(x_k\in[a,b]\) converge to \(x\). It
suffices to show \(g(x_k)\to g(x)\). Fix \(y\neq x\). If \(y<a\), then
\(\mathbf1_{[a,x_k]}(y)=0=\mathbf1_{[a,x]}(y)\) for all \(k\). If
\(a\leq y<x\), then \(x_k>y\) for all large \(k\), so
\(\mathbf1_{[a,x_k]}(y)=1=\mathbf1_{[a,x]}(y)\) eventually. If
\(y>x\), then \(x_k<y\) for all large \(k\), so both indicators are
eventually zero. Hence
\(f\mathbf1_{[a,x_k]}\to f\mathbf1_{[a,x]}\) at every point except
possibly \(x\), and \(\lambda(\{x\})=0\) by
\cref{thm:p15-lebesgue-unit-interval}. Every
\(|f\mathbf1_{[a,x_k]}|\) is at most \(|f|\mathbf1_{[a,b]}\), which is
integrable by hypothesis. Dominated convergence
(\cref{thm:p15-dominated-bounded-convergence}~(1)) gives
\(g(x_k)=\int f\mathbf1_{[a,x_k]}\,d\lambda\to
\int f\mathbf1_{[a,x]}\,d\lambda=g(x)\). Each \(g(x)\) is finite
because \(|f\mathbf1_{[a,x]}|\leq|f|\mathbf1_{[a,b]}\). So \(g\) is
continuous at every point of \([a,b]\), in particular on \((a,b)\).
\end{proof}
\end{annexitem}

\begin{annexitem}{annex:p15-w02c-p15-005-ae-measure}{thm:p15-w02c-p15-005-ae-measure}{Continuity from above on a finite measure space}
\begin{proof}
(1) Let \(N\) be a measurable null set off which \(f_n\to f\), and fix
\(\varepsilon>0\). Put
\[
 B_n\defeq\bigcup_{m\geq n}\{|f_m-f|>\varepsilon\}.
\]
These sets are measurable and decrease in \(n\). If
\(\omega\in\bigcap_nB_n\), then \(|f_m(\omega)-f(\omega)|>\varepsilon\)
for infinitely many \(m\), so \(f_m(\omega)\not\to f(\omega)\) and
\(\omega\in N\). Since \(\mu(B_1)\leq\mu(\Omega)<\infty\), continuity
of the measure from above applies:
\(\mu(B_n)\to\mu\bigl(\bigcap_nB_n\bigr)\leq\mu(N)=0\). As
\(\{|f_n-f|>\varepsilon\}\subseteq B_n\), it follows that
\(\mu(\{|f_n-f|>\varepsilon\})\to0\). This holds for every
\(\varepsilon>0\).

(2) For each real \(y\), \(\mathbf1_{(n,n+1]}(y)=0\) as soon as
\(n\geq y\), so \(f_n\to0\) everywhere. For every \(n\),
\(\{|f_n|>1/2\}=(n,n+1]\), of Lebesgue measure one by
\cref{thm:p15-lebesgue-unit-interval}. So \(f_n\) does not converge to
zero in measure; the only finite-measure step in part~(1), continuity
from above, fails here because the sets
\(B_n=(n,\infty)\) have infinite measure and empty intersection.

(3) By \cref{thm:p15-w02b-p15-013-countermodels}~(2), the typewriter
sequence has \(\mathbb E|T_n|=2^{-k}\to0\), so
\(\mathbb U(|T_n|>\varepsilon)\leq\mathbb E|T_n|/\varepsilon\to0\)
for every \(\varepsilon>0\) by Markov's bound
(\cref{thm:p15-expectation-bounds}), and at every point of \((0,1)\)
it takes the values zero and one infinitely often, so it converges at
no point.
\end{proof}
\end{annexitem}

\begin{annexitem}{annex:p15-w02c-p15-005-ae-bounded}{thm:p15-w02c-p15-005-ae-bounded}{Continuity from above inside a set of finite measure}
\begin{proof}
Each hypothesis holds off a measurable null set, and there are
countably many of them, so there is a measurable null set \(N\) such
that, for \(\omega\notin N\), \(f_n(\omega)\to f(\omega)\),
\(|f_n(\omega)|\leq M\) for every \(n\), and \(f_n(\omega)=0\) for
every \(n\) when \(\omega\notin E\). For \(\omega\notin E\cup N\) it
follows that \(f(\omega)=\lim_nf_n(\omega)=0\), so
\(|f_n(\omega)-f(\omega)|=0\) for every \(n\).

Fix \(\varepsilon>0\) and put
\(B_n\defeq E\cap\bigcup_{m\geq n}\{|f_m-f|>\varepsilon\}\). These
measurable sets decrease in \(n\), and \(\mu(B_1)\leq\mu(E)<\infty\).
A point of \(\bigcap_nB_n\) has \(|f_m-f|>\varepsilon\) for infinitely
many \(m\), so \(f_m\) does not converge to \(f\) there, and the point
lies in \(N\). Continuity of \(\mu\) from above gives
\(\mu(B_n)\to\mu\bigl(\bigcap_nB_n\bigr)\leq\mu(N)=0\). By the first
paragraph,
\(\{|f_n-f|>\varepsilon\}\subseteq B_n\cup N\), so
\(\mu(\{|f_n-f|>\varepsilon\})\leq\mu(B_n)\to0\). As \(\varepsilon\)
was arbitrary, \(f_n\to f\) in measure in the sense of
\cref{def:p15-lp-convergence}.

Now every hypothesis of \cref{thm:p15-dominated-bounded-convergence}~(2)
holds: \(\mu(E)<\infty\), each \(f_n\) vanishes outside \(E\) a.e.,
\(|f_n|\leq M\) a.e., and \(f_n\to f\) in measure. That theorem gives
\(\|f_n-f\|_1\to0\) and \(\int f_n\,d\mu\to\int f\,d\mu\).
\end{proof}
\end{annexitem}

\begin{annexitem}{annex:p15-w02c-p15-005-supporting-affine}{thm:p15-w02c-p15-005-supporting-affine}{A cube bound and a linear minorant}
\begin{proof}
\emph{Continuity.} Fix \(c\in\mathbb R^n\) and let
\(Q=c+[-1,1]^n\). Every point of \(Q\) is a convex combination of the
\(2^n\) vertices \(c+\epsilon\), \(\epsilon\in\{-1,1\}^n\): writing
\(x_i=c_i+t_i-(1-t_i)\) with \(t_i\in[0,1]\), the weights
\(\prod_i(t_i\text{ or }1-t_i)\), with \(t_i\) chosen when
\(\epsilon_i=1\), are nonnegative, sum to one, and reproduce \(x\)
coordinate by coordinate. The convexity inequality extends to finite
convex combinations by induction on the number \(k\) of points: if
\(w_1<1\), then
\(\sum_{j\leq k}w_jz_j=w_1z_1+(1-w_1)\sum_{2\leq j\leq k}\frac{w_j}{1-w_1}z_j\),
and the inner combination has \(k-1\) points; if \(w_1=1\), all other
weights vanish. Hence \(\varphi\leq M\) on \(Q\), where \(M\) is the
largest value of \(\varphi\) at a vertex. Now let \(x\in Q\),
\(x\neq c\), and \(r=\max_i|x_i-c_i|\in(0,1]\). Put
\(u=(x-c)/r\), so \(c+u\) and \(c-u\) lie in \(Q\). Then
\(x=(1-r)c+r(c+u)\) and \(c=\frac1{1+r}x+\frac r{1+r}(c-u)\). Convexity
gives
\[
 \varphi(x)\leq(1-r)\varphi(c)+rM,\qquad
 (1+r)\varphi(c)\leq\varphi(x)+rM .
\]
Hence \(|\varphi(x)-\varphi(c)|\leq r\,(M-\varphi(c))\), which tends to
zero as \(x\to c\). So \(\varphi\) is continuous. Each set
\(\{\varphi>t\}\) is then open, hence Borel, so \(\varphi\) is Borel
measurable.

\emph{A sublinear minorant.} For \(v\in\mathbb R^n\) and \(s>0\), put
\(q_v(s)=(\varphi(c+sv)-\varphi(c))/s\). For \(0<s<s'\), writing
\(c+sv=(1-s/s')c+(s/s')(c+s'v)\) and using convexity gives
\(q_v(s)\leq q_v(s')\). Writing
\(c=\frac{s'}{s+s'}(c+sv)+\frac s{s+s'}(c-s'v)\) gives
\(q_v(s)\geq(\varphi(c)-\varphi(c-s'v))/s'\). So \(q_v\) is
nondecreasing and bounded below, and
\[
 p(v)\defeq\lim_{s\downarrow0}q_v(s)=\inf_{s>0}q_v(s)
\]
is a real number with \(p(v)\leq q_v(1)=\varphi(c+v)-\varphi(c)\). For
\(\alpha>0\), \(q_{\alpha v}(s)=\alpha q_v(\alpha s)\), so
\(p(\alpha v)=\alpha p(v)\); and \(p(0)=0\). For \(u,v\), convexity at
the midpoint gives
\(\varphi(c+s(u+v))\leq\frac12\varphi(c+2su)+\frac12\varphi(c+2sv)\),
that is \(q_{u+v}(s)\leq q_u(2s)+q_v(2s)\), and letting
\(s\downarrow0\) gives \(p(u+v)\leq p(u)+p(v)\).

\emph{A linear functional below \(p\).} We show that for each linear
subspace \(V\subseteq\mathbb R^n\) with a linear \(L:V\to\mathbb R\),
\(L\leq p\) on \(V\), and each \(w\notin V\), the map \(L\) extends
linearly to \(V+\mathbb Rw\), still below \(p\). For \(x,y\in V\),
\[
 L(x)+L(y)=L(x+y)\leq p(x+y)\leq p(x+w)+p(y-w),
\]
so \(L(y)-p(y-w)\leq p(x+w)-L(x)\). Hence
\(\sup_{y\in V}(L(y)-p(y-w))\leq\inf_{x\in V}(p(x+w)-L(x))\), both
sides being real since \(p\) is real valued; choose \(\beta\) between
them and set \(L(x+tw)=L(x)+t\beta\). For \(t>0\),
\(L(x)+t\beta=t\,(L(x/t)+\beta)\leq t\,p(x/t+w)=p(x+tw)\). For
\(t=-s<0\), \(L(x)-s\beta=s\,(L(x/s)-\beta)\leq s\,p(x/s-w)=p(x-sw)\).
Starting from \(V=\{0\}\), \(L=0\), and adjoining the standard basis
vectors not yet in the span one at a time, after at most \(n\) steps
we obtain a linear \(L\) on \(\mathbb R^n\) with \(L\leq p\). Write
\(L(v)=m\cdot v\) with \(m_i=L(e_i)\).

\emph{The supporting inequality.} For every \(x\), with \(v=x-c\),
\[
 \varphi(x)-\varphi(c)=q_v(1)\geq p(v)\geq m\cdot(x-c).
\]

\emph{Strictness.} Let \(\varphi\) be strictly convex and suppose
\(\varphi(x)=\varphi(c)+m\cdot(x-c)\) for some \(x\neq c\). With
\(z=\frac12(x+c)\), strict convexity gives
\(\varphi(z)<\frac12\varphi(x)+\frac12\varphi(c)
=\varphi(c)+m\cdot(z-c)\), contradicting the supporting inequality at
\(z\). So the inequality is strict for every \(x\neq c\).
\end{proof}
\end{annexitem}

\begin{annexitem}{annex:p15-w02c-p15-005-jensen-several}{thm:p15-w02c-p15-005-jensen-several}{Integrating the supporting affine function}
\begin{proof}
Write \(X=(X_1,\ldots,X_n)\). For an open set \(O\subseteq\mathbb R^n\),
\(O\) is the union of the countably many open boxes
\(\prod_i(a_i,b_i)\) with rational endpoints contained in \(O\), and
\(\{X\in\prod_i(a_i,b_i)\}=\bigcap_i\{a_i<X_i<b_i\}\) is an event. So
\(\{X\in O\}\) is an event. By
\cref{thm:p15-w02c-p15-005-supporting-affine}, \(\varphi\) is
continuous, and \(\{\varphi(X)>t\}=\{X\in\{\varphi>t\}\}\) is an event
for every real \(t\); hence \(\varphi(X)\) is a random variable.

Let \(c=(\mathbb EX_1,\ldots,\mathbb EX_n)\), a point of
\(\mathbb R^n\) because each \(X_i\) is integrable, and take \(m\) from
\cref{thm:p15-w02c-p15-005-supporting-affine}. Then pointwise
\[
 \varphi(X)\geq\varphi(c)+\sum_{i=1}^nm_i(X_i-c_i).
\]
The right side is integrable, a finite linear combination of
integrable variables and a constant. Both sides are integrable, so
taking expectations with linearity and order
(\cref{thm:p15-signed-integral-laws}) gives
\[
 \mathbb E\varphi(X)\geq\varphi(c)+\sum_{i=1}^nm_i(\mathbb EX_i-c_i)
 =\varphi(c)=\varphi(\mathbb EX_1,\ldots,\mathbb EX_n).
\]
\end{proof}
\end{annexitem}

\begin{annexitem}{annex:p15-w02c-p15-005-jensen-special}{thm:p15-w02c-p15-005-jensen-special}{Jensen with the absolute value and the square}
\begin{proof}
The function \(x\mapsto|x|\) is convex, by the triangle inequality:
\(|tx+(1-t)y|\leq t|x|+(1-t)|y|\). With \(\mathbb E|X|<\infty\), the
variable \(|X|\) is integrable, and Jensen's inequality
(\cref{thm:p15-jensen-integral}) gives
\(|\mathbb EX|\leq\mathbb E|X|\).

The function \(x\mapsto x^2\) is convex, since
\(tx^2+(1-t)y^2-(tx+(1-t)y)^2=t(1-t)(x-y)^2\geq0\). If
\(\mathbb E(X^2)=\infty\), the second inequality holds because its
left side is a finite real number. If \(\mathbb E(X^2)<\infty\), then
\(X^2\) is integrable, and Jensen's inequality gives
\((\mathbb EX)^2\leq\mathbb E(X^2)\).
\end{proof}
\end{annexitem}

\begin{annexitem}{annex:p15-w02c-p15-005-jensen-equality}{thm:p15-w02c-p15-005-jensen-equality}{A nonnegative variable with zero mean}
\begin{proof}
Let \(c=\mathbb EX\). If \(X=c\) a.s., then \(\varphi(X)=\varphi(c)\)
a.s., so \(\mathbb E\varphi(X)=\varphi(c)\).

Conversely, suppose \(\varphi(c)=\mathbb E\varphi(X)\). By
\cref{thm:p15-w02c-p15-005-supporting-affine} with \(n=1\), there is a
real \(m\) with \(\ell(x)\defeq\varphi(c)+m(x-c)\leq\varphi(x)\) for
all \(x\), with strict inequality for \(x\neq c\). The variable
\(D=\varphi(X)-\ell(X)\) is nonnegative and integrable, and by linearity
\[
 \mathbb ED=\mathbb E\varphi(X)-\varphi(c)-m(\mathbb EX-c)=0 .
\]
By \cref{thm:p15-nonnegative-integral-laws}~(1), \(D=0\) a.s. Since
\(D(\omega)>0\) whenever \(X(\omega)\neq c\), the event
\(\{X\neq c\}\) is contained in \(\{D>0\}\), which is null. So
\(X=c\) a.s.
\end{proof}
\end{annexitem}

\begin{annexitem}{annex:p15-w02c-p15-005-am-gm}{thm:p15-w02c-p15-005-am-gm}{Jensen for the exponential on a finite space}
\begin{proof}
\emph{The exponential is strictly convex.} Let
\(k(z)=e^z-1-z\). It is differentiable with \(k'(z)=e^z-1\), which is
negative for \(z<0\) and positive for \(z>0\). By the mean value
theorem \(k\) is strictly decreasing on \((-\infty,0]\) and strictly
increasing on \([0,\infty)\), so \(k(z)>k(0)=0\) for \(z\neq0\); that
is, \(e^z>1+z\) for \(z\neq0\). Let \(x\neq y\), \(0<t<1\), and
\(w=tx+(1-t)y\); then \(x\neq w\) and \(y\neq w\). So
\(e^x=e^we^{x-w}>e^w(1+x-w)\) and \(e^y>e^w(1+y-w)\). Combining with
weights \(t\) and \(1-t\),
\[
 te^x+(1-t)e^y>e^w\bigl(1+t(x-w)+(1-t)(y-w)\bigr)=e^w .
\]

\emph{The inequality.} Let \(\Omega=\{1,\ldots,n\}\) with all subsets as
events and \(\mathbb P(\{m\})=p(m)\), a probability since the weights are
positive with sum one, and let \(X(m)=\log y_m\). Then
\(\mathbb P(X=\log y_m)\geq p(m)\), with equality when the \(y_m\) are
distinct; in every case expectations are the weighted sums
\(\sum_mp(m)(\cdot)\) of \cref{thm:p15-expectation-law-integral}. Both
\(X\) and \(e^X\) are integrable on this finite space, and
\[
 \mathbb Ee^X=\sum_{m=1}^np(m)y_m,
 \qquad
 e^{\mathbb EX}=\exp\Bigl(\sum_{m=1}^np(m)\log y_m\Bigr)
 =\prod_{m=1}^ny_m^{p(m)} .
\]
Jensen's inequality (\cref{thm:p15-jensen-integral}) gives the
displayed inequality.

\emph{Equality.} If all \(y_m\) are equal to \(y\), both sides equal
\(y\). Conversely, if equality holds, then by
\cref{thm:p15-w02c-p15-005-jensen-equality} \(X=\mathbb EX\) a.s.;
since every point has positive probability, \(X\) is constant on
\(\Omega\), so \(\log y_1=\cdots=\log y_n\) and the \(y_m\) are equal.
With \(p(m)=1/n\) the inequality reads
\(\frac1n\sum_my_m\geq\bigl(\prod_my_m\bigr)^{1/n}\).
\end{proof}
\end{annexitem}

\begin{annexitem}{annex:p15-w02c-p15-005-chebyshev-displays}{thm:p15-w02c-p15-005-chebyshev-displays}{The variance identity}
\begin{proof}
Since \(\mathbb EX^2<\infty\), \cref{thm:p15-expectation-bounds} shows
that \(X\) is integrable and
\(\operatorname{Var}(X)=\mathbb EX^2-(\mathbb EX)^2\), a difference of
two finite real numbers. As \((\mathbb EX)^2\geq0\),
\(\operatorname{Var}(X)\leq\mathbb EX^2\), and the two sides are equal
exactly when \((\mathbb EX)^2=0\), that is, when \(\mathbb EX=0\).
\end{proof}
\end{annexitem}

\begin{annexitem}{annex:p15-w02c-p15-005-chebyshev-sharp}{thm:p15-w02c-p15-005-chebyshev-sharp}{A two-point extremal law and a vanishing tail second moment}
\begin{proof}
(1) Since \(0<b\leq a\), \(0<b^2/a^2\leq1\), so the two masses are a
probability law; it is realized, for example, on a two-point space.
Then \(\mathbb EX^2=a^2\cdot b^2/a^2+0=b^2\), and since \(a>0\),
\(\{|X|\geq a\}=\{X=a\}\), of probability \(b^2/a^2\). Thus
\(a^2\mathbb P(|X|\geq a)=b^2=\mathbb EX^2\).

(2) For \(a>0\), apply \cref{thm:p15-expectation-bounds} with the
nonnegative Borel function \(\varphi(x)=x^2\) and the nonempty Borel
set \(A=\{x:|x|\geq a\}\), on which the infimum of \(\varphi\) is
\(a^2\), attained at \(x=a\). It gives
\(a^2\mathbb P(|X|\geq a)\leq\mathbb E(X^2;|X|\geq a)\). The function
\(a\mapsto\mathbb E(X^2;|X|\geq a)\) is nonincreasing, so it suffices to
show that it tends to zero along \(a=k\), \(k=1,2,\ldots\). Since \(X\)
is real valued, \(X^2\mathbf1_{\{|X|\geq k\}}\to0\) at every point, and
\(X^2\mathbf1_{\{|X|\geq k\}}\leq X^2\), which is integrable. Dominated
convergence (\cref{thm:p15-dominated-bounded-convergence}~(1)) gives
\(\mathbb E(X^2;|X|\geq k)\to0\). Dividing by the positive finite
number \(\mathbb EX^2\) gives the limit zero.
\end{proof}
\end{annexitem}

\begin{annexitem}{annex:p15-w02c-p15-005-one-sided}{thm:p15-w02c-p15-005-one-sided}{Chebyshev with a shifted square}
\begin{proof}
(1) Let \(\varphi(y)=(y+b)^2\) and \(A=[a,\infty)\). For \(y\geq a\),
\(y+b\geq a+b>0\), so \(\inf_A\varphi=(a+b)^2\).
\Cref{thm:p15-expectation-bounds} gives
\((a+b)^2\mathbb P(Y\geq a)\leq\mathbb E(Y+b)^2\). By linearity and the
variance identity,
\[
 \mathbb E(Y+b)^2=\mathbb EY^2+2b\,\mathbb EY+b^2
 =\operatorname{Var}(Y)+(\mathbb EY+b)^2 .
\]
The same identity holds for \(X\), whose second moment is finite. Since
\(Y\) and \(X\) have the same mean and variance,
\[
 \mathbb E(Y+b)^2=\mathbb E(X+b)^2=p\,(a+b)^2+(1-p)\cdot0 .
\]
Dividing by \((a+b)^2>0\) gives \(\mathbb P(Y\geq a)\leq p\). For
\(Y=X\), \(\{X\geq a\}=\{X=a\}\) because \(-b<0<a\), so
\(\mathbb P(X\geq a)=p\).

(2) If \(\sigma=0\), then \(\mathbb EY^2=0\), so \(Y=0\) a.s. and
\(\mathbb P(Y\geq a)=0=\sigma^2/(a^2+\sigma^2)\); the variable \(Y=0\)
attains equality. Let \(\sigma>0\), and put
\[
 b=\frac{\sigma^2}a,\qquad p=\frac{\sigma^2}{a^2+\sigma^2}.
\]
Then \(b>0\) and \(0<p<1\). For the law \(\mathbb P(X=a)=p\),
\(\mathbb P(X=-b)=1-p\),
\[
 \mathbb EX=\frac{a\sigma^2-\frac{\sigma^2}a\,a^2}{a^2+\sigma^2}=0,
 \qquad
 \mathbb EX^2=\frac{\sigma^2a^2+\frac{\sigma^4}{a^2}a^2}{a^2+\sigma^2}
 =\sigma^2 .
\]
So \(X\) has mean zero and variance \(\sigma^2\), the same as \(Y\).
Part~(1) gives \(\mathbb P(Y\geq a)\leq p=\sigma^2/(a^2+\sigma^2)\), with
equality for \(Y=X\).
\end{proof}
\end{annexitem}

\begin{annexitem}{annex:p15-w02c-p15-005-no-lower}{thm:p15-w02c-p15-005-no-lower}{Two-point laws with a rare far value}
\begin{proof}
Every infimum of probabilities is at least zero, so it suffices in each
part to exhibit admissible variables whose probabilities tend to zero.
A law with finitely many atoms is realized on a finite space, and its
moments are the weighted sums of \cref{thm:p15-expectation-law-integral}.

(1) For an integer \(n\geq1\), let
\(\mathbb P(X_n=n)=\mathbb P(X_n=-n)=1/(2n^2)\) and
\(\mathbb P(X_n=0)=1-1/n^2\). Then \(\mathbb EX_n=0\) and
\(\operatorname{Var}(X_n)=\mathbb EX_n^2=n^2\cdot1/n^2=1\). For
\(n>\varepsilon\), \(\mathbb P(|X_n|>\varepsilon)=1/n^2\to0\).

(2) For \(0<r<1\), let a two-point law put mass \(1-r\) at
\(u_r\) and mass \(r\) at \(v_r\), where
\[
 d_r=\frac{\sigma}{\sqrt{r(1-r)}},\qquad u_r=1-rd_r,\qquad v_r=u_r+d_r .
\]
Its mean is \((1-r)u_r+r(u_r+d_r)=u_r+rd_r=1\). Its variance is
\[
 (1-r)(u_r-1)^2+r(v_r-1)^2=(1-r)r^2d_r^2+r(1-r)^2d_r^2
 =r(1-r)d_r^2=\sigma^2 .
\]
Now \(rd_r=\sigma\sqrt{r/(1-r)}\to0\) as \(r\downarrow0\), so for all
small \(r\), \(0\leq rd_r\leq1+y\), whence
\(-y\leq u_r\leq1\leq y\) and \(|u_r|\leq y\). For such \(r\), only the
atom \(v_r\) can satisfy \(|x|>y\), so
\(\mathbb P(|X|>y)\leq r\to0\).
\end{proof}
\end{annexitem}

\begin{annexitem}{annex:p15-w02c-p15-005-positive-lower}{thm:p15-w02c-p15-005-positive-lower}{Cauchy--Schwarz on the positive region}
\begin{proof}
Since \(Y\geq0\), \(Y=Y\mathbf1_{(Y>0)}\). Hölder's inequality with
\(p=q=2\) (\cref{thm:p15-holder-integral}) applied to \(Y\) and
\(\mathbf1_{(Y>0)}\) gives
\[
 \mathbb EY=\mathbb E\bigl(Y\mathbf1_{(Y>0)}\bigr)
 \leq(\mathbb EY^2)^{1/2}\bigl(\mathbb E\mathbf1_{(Y>0)}\bigr)^{1/2}
 =(\mathbb EY^2)^{1/2}\,\mathbb P(Y>0)^{1/2}.
\]
Both sides are finite and nonnegative; squaring gives
\((\mathbb EY)^2\leq\mathbb EY^2\cdot\mathbb P(Y>0)\). If
\(\mathbb EY^2>0\), divide by it. If \(\mathbb EY^2=0\), then \(Y^2=0\)
a.s. by \cref{thm:p15-nonnegative-integral-laws}~(1), so \(Y=0\) a.s.
and \(\mathbb P(Y>0)=0\).
\end{proof}
\end{annexitem}

\begin{annexitem}{annex:p15-w02c-p15-005-moment-power}{thm:p15-w02c-p15-005-moment-power}{The power case of moment-controlled convergence}
\begin{proof}
Put \(g(x)=|x|^p\) and \(h(x)=x\) for real \(x\). Both are continuous,
and \(g\geq0\). Since \(p>0\), \(g(x)\to\infty\) as \(|x|\to\infty\).
For \(x\neq0\), \(g(x)>0\) and
\(|h(x)|/g(x)=|x|^{1-p}\), which tends to \(0\) as \(|x|\to\infty\)
because \(1-p<0\). By hypothesis \(\mathbb Eg(X_n)\leq K<\infty\) for
every \(n\), and \(X_n\to X\) a.s. Every hypothesis of
\cref{thm:p15-moment-controlled-convergence} therefore holds, and that
theorem gives that \(h(X_n)=X_n\) and \(h(X)=X\) are integrable and
\(\mathbb EX_n\to\mathbb EX\).

For the moment bound on the limit, \(g\) is continuous, so
\(g(X_n)\to g(X)\) a.s., and hence \(\liminf_ng(X_n)=g(X)\) a.s. The
functions \(g(X_n)\) are nonnegative and measurable, so Fatou's
inequality, \cref{thm:p15-nonnegative-integral-laws}~(3), together with
part~(2) of that theorem for the a.s.\ equality, gives
\[
 \mathbb E|X|^p=\mathbb E\liminf_ng(X_n)\leq\liminf_n\mathbb Eg(X_n)
 \leq K .
\]

For the remark after the corollary: \(\mathbb U((0,1/n))=1/n\) by
\cref{thm:p15-lebesgue-unit-interval}, so
\(\mathbb E|Y_n|=\mathbb EY_n=n\cdot(1/n)=1\). A point
\(u\in(0,1)\) lies outside \((0,1/n)\) once \(n>1/u\), so
\(Y_n(u)=0\) for all large \(n\) and \(Y_n\to0\) at every point, while
\(\mathbb E0=0\neq1\).
\end{proof}
\end{annexitem}

\begin{annexitem}{annex:p15-w02c-p15-005-no-dominator}{thm:p15-w02c-p15-005-no-dominator}{A spike sequence below a fractional power}
\begin{proof}
Write \(I_n=(1/(n+1),1/n)\); these intervals are disjoint, and
\(\mathbb U(I_n)=1/n-1/(n+1)=1/(n(n+1))\) by
\cref{thm:p15-lebesgue-unit-interval}.

(1) A point \(u\in(0,1)\) lies in at most one \(I_n\), so
\(X_n(u)=0\) for all but at most one \(n\), and \(X_n(u)\to0\).

(2) The functions \(g\) and \(h\) are continuous, \(g\geq0\), and
\(g(x)\to\infty\) as \(|x|\to\infty\) because \(2/\alpha>0\). For
\(x\neq0\), \(|h(x)|/g(x)=|x|^{1-2/\alpha}\to0\) as \(|x|\to\infty\)
because \(1-2/\alpha<0\) when \(\alpha<2\). Since \(g(0)=0\),
\[
 \mathbb Eg(X_n)=(n^\alpha)^{2/\alpha}\,\mathbb U(I_n)
 =\frac{n^2}{n(n+1)}=\frac n{n+1}\leq1 .
\]
So \cref{thm:p15-moment-controlled-convergence} applies with \(K=1\)
and \(X=0\), and gives \(\mathbb EX_n\to0\). Directly,
\(\mathbb EX_n=n^\alpha/(n(n+1))=n^{\alpha-1}/(n+1)\leq n^{\alpha-2}\to0\)
because \(\alpha<2\).

(3) Suppose \(Y\) is a random variable with \(|X_n|\leq Y\) a.s.\ for
every \(n\). Then \(Y\geq n^\alpha\) a.s.\ on \(I_n\), and since the
\(I_n\) are disjoint,
\[
 \mathbb E|Y|\geq\sum_{n\geq1}n^\alpha\,\mathbb U(I_n)
 =\sum_{n\geq1}\frac{n^\alpha}{n(n+1)}
 \geq\sum_{n\geq1}\frac{n^{\alpha-2}}2
 \geq\frac12\sum_{n\geq1}\frac1n ,
\]
using \(n+1\leq2n\) and \(n^{\alpha-2}\geq n^{-1}\) for \(n\geq1\),
which holds because \(\alpha-2>-1\). The first inequality uses
countable additivity of \(A\mapsto\mathbb E(|Y|;A)\)
(\cref{thm:p15-nonnegative-integral-laws}~(4)). The harmonic series
diverges: for \(k\geq0\), the \(2^k\) terms with
\(2^k<n\leq2^{k+1}\) are each at least \(2^{-k-1}\), so they sum to at
least \(1/2\), and there are infinitely many such blocks. Hence
\(\mathbb E|Y|=\infty\), and no integrable dominator exists.
\end{proof}
\end{annexitem}

\begin{annexitem}{annex:p15-w02c-p15-005-inverse-tail}{thm:p15-w02c-p15-005-inverse-tail}{A bounded ratio and bounded convergence}
\begin{proof}
For \(y>0\), let \(Z_y=(y/X)\mathbf1_{(X>y)}\), with value zero off
\(\{X>y\}\). On \(\{X>y\}\), \(0<y/X<1\), so \(0\leq Z_y\leq1\) and
\(y\,\mathbb E(1/X;X>y)=\mathbb EZ_y\leq\mathbb P(X>y)\leq1\); dividing by
\(y\) gives \(\mathbb E(1/X;X>y)\leq1/y\).

As \(y\to\infty\): \(\mathbb EZ_y\leq\mathbb P(X>y)\). The events
\(\{X>k\}\), \(k=1,2,\ldots\), decrease to the empty set because \(X\)
is real valued, so \(\mathbb P(X>k)\to0\) by continuity from above;
since \(\mathbb P(X>y)\) is nonincreasing in \(y\),
\(\mathbb P(X>y)\to0\) as \(y\to\infty\), and so does
\(\mathbb EZ_y\).

As \(y\downarrow0\): let \(y_k\downarrow0\) be any sequence of positive
numbers. At a point where \(X=0\), \(Z_{y_k}=0\) for every \(k\). At a
point where \(X>0\), \(0\leq Z_{y_k}\leq y_k/X\to0\). So
\(Z_{y_k}\to0\) everywhere, with \(|Z_{y_k}|\leq1\), and the constant one
is integrable on a probability space. Dominated convergence
(\cref{thm:p15-dominated-bounded-convergence}~(1)) gives
\(\mathbb EZ_{y_k}\to0\). As this holds for every sequence decreasing to
zero, \(\lim_{y\downarrow0}\mathbb EZ_y=0\): otherwise some
\(\eta>0\) and positive \(y_k\to0\) would have
\(\mathbb EZ_{y_k}\geq\eta\), and a decreasing subsequence of
\((y_k)\) would contradict what was just shown.
\end{proof}
\end{annexitem}

\begin{annexitem}{annex:p15-w02c-p15-005-inclusion-exclusion}{thm:p15-w02c-p15-005-inclusion-exclusion}{Expanding a product of complementary indicators}
\begin{proof}
At a point of \(A\), some \(\mathbf1_{A_i}\) equals one, so the product
\(\prod_i(1-\mathbf1_{A_i})\) is zero and
\(1-\prod_i(1-\mathbf1_{A_i})=1=\mathbf1_A\). At a point outside \(A\),
every factor is one, and both sides are zero. Expanding the product of
the \(n\) binomials, choosing from each factor either \(1\) or
\(-\mathbf1_{A_i}\), and using
\(\prod_{i\in S}\mathbf1_{A_i}=\mathbf1_{A_S}\),
\[
 \prod_{i=1}^n(1-\mathbf1_{A_i})
 =\sum_{S\subseteq\{1,\ldots,n\}}(-1)^{|S|}\mathbf1_{A_S},
\]
where the term for the empty set \(S\) is the constant one. Hence
\(\mathbf1_A=\sum_{S\neq\varnothing}(-1)^{|S|-1}\mathbf1_{A_S}\). Taking
expectations, with linearity of expectation for this finite sum of
bounded variables and \(\mathbb E\mathbf1_B=\mathbb P(B)\)
(\cref{thm:p15-expectation-law-integral}), and grouping the sets \(S\)
by their size \(k\), gives
\(\mathbb P(A)=\sum_{k=1}^n(-1)^{k-1}S_k\).
\end{proof}
\end{annexitem}

\begin{annexitem}{annex:p15-w02c-p15-005-bonferroni}{thm:p15-w02c-p15-005-bonferroni}{Counting the events that occur}
\begin{proof}
For a nonnegative integer \(N\) and an integer \(m\geq0\), write
\(\binom Nm\) for the number of \(m\)-element subsets of an
\(N\)-element set, zero when \(m>N\). We first show, for \(N\geq1\)
and \(m\geq0\),
\[
 \sum_{k=0}^m(-1)^k\binom Nk=(-1)^m\binom{N-1}m .
\]
For \(m=0\) both sides are one. If it holds for \(m\), then by
Pascal's rule \(\binom N{m+1}=\binom{N-1}{m+1}+\binom{N-1}m\),
\[
 \begin{aligned}
 \sum_{k=0}^{m+1}(-1)^k\binom Nk
 &=(-1)^m\binom{N-1}m+(-1)^{m+1}\binom{N-1}m\\
 &\qquad+(-1)^{m+1}\binom{N-1}{m+1}\\
 &=(-1)^{m+1}\binom{N-1}{m+1}.
 \end{aligned}
\]
Pascal's rule holds because the \((m+1)\)-subsets of an \(N\)-set either
avoid or contain a fixed element, and these are counted by
\(\binom{N-1}{m+1}\) and \(\binom{N-1}m\).

Fix \(\omega\) and let \(N(\omega)\) be the number of indices \(i\) with
\(\omega\in A_i\). A set \(S\) of size \(k\) has \(\omega\in A_S\)
exactly when \(S\) is a subset of those \(N(\omega)\) indices, so
\(\sum_{|S|=k}\mathbf1_{A_S}(\omega)=\binom{N(\omega)}k\). Let
\(T_m=\sum_{k=1}^m(-1)^{k-1}\sum_{|S|=k}\mathbf1_{A_S}\). If
\(N(\omega)=0\), then \(T_m(\omega)=0=\mathbf1_A(\omega)\). If
\(N(\omega)\geq1\), then \(\mathbf1_A(\omega)=1\) and, by the identity,
\[
 T_m(\omega)=1-\sum_{k=0}^m(-1)^k\binom{N(\omega)}k
 =\mathbf1_A(\omega)+(-1)^{m+1}\binom{N(\omega)-1}m .
\]
So \(T_m-\mathbf1_A\geq0\) everywhere when \(m\) is odd and
\(T_m-\mathbf1_A\leq0\) everywhere when \(m\) is even. For \(m=1\) this
is \(\mathbf1_A\leq\sum_i\mathbf1_{A_i}\). Taking expectations, which
preserves order (\cref{thm:p15-signed-integral-laws}), and using
\(\mathbb ET_m=\sum_{k=1}^m(-1)^{k-1}S_k\) by linearity, gives the two
inequalities. For \(m=n\), \(\binom{N(\omega)-1}n=0\) because
\(N(\omega)-1<n\), so \(T_n=\mathbf1_A\), which recovers
\cref{thm:p15-w02c-p15-005-inclusion-exclusion}.
\end{proof}
\end{annexitem}
